% Vietnamese translation for OI Public Library
% Source-PDF: 英文资料 ENGLISH MATERIALS/Data Structures and Network Algorithms - Tarjan [1987-01-01].pdf
\documentclass[11pt]{article}
\usepackage{oipl-vn}

\newcommand{\Oh}{\mathcal{O}}

\title{Cấu trúc dữ liệu và thuật toán mạng}
\author{Robert Endre Tarjan}
\date{SIAM, 1983}

\begin{document}
\maketitle

\origtitle{Data Structures and Network Algorithms}
\sourcepath{English Materials / Data Structures and Network Algorithms - Tarjan [1987-01-01].pdf}

\renewcommand{\abstractname}{Tóm tắt}
\begin{abstract}
Đây là ghi chú dịch tiếng Việt cho chuyên khảo ngắn của Robert E. Tarjan về
mối liên hệ giữa cấu trúc dữ liệu, phân tích thuật toán và tối ưu mạng. Bản
dịch trong kho này chuyển ngữ phần thông tin đầu sách, mục lục và lời nói đầu,
đồng thời ghi lại nội dung chính của từng chương để phục vụ tra cứu trong bối
cảnh luyện thi thuật toán.
\end{abstract}

\section*{Thông tin sách}

Sách thuộc bộ \emph{CBMS-NSF Regional Conference Series in Applied
Mathematics}, do SIAM xuất bản. Bản gốc:

\begin{center}
\begin{tabular}{ll}
\textbf{Tác giả} & Robert Endre Tarjan \\
\textbf{Cơ quan} & Bell Laboratories, Murray Hill, New Jersey \\
\textbf{Nhà xuất bản} & Society for Industrial and Applied Mathematics \\
\textbf{Năm} & 1983 \\
\textbf{ISBN} & 0-89871-187-8
\end{tabular}
\end{center}

\section{Mục lục đã dịch}

\begin{description}
  \item[Chương 1. Nền tảng]
  Giới thiệu; độ phức tạp tính toán; cấu trúc dữ liệu nguyên thủy; ký hiệu
  thuật toán; cây và đồ thị.

  \item[Chương 2. Các tập rời nhau]
  Tập rời nhau và cây nén; cận trên khấu hao cho nén đường đi; ghi chú.

  \item[Chương 3. Heap]
  Heap và cây có thứ tự heap; ``cheap''; heap trái; ghi chú.

  \item[Chương 4. Cây tìm kiếm]
  Tập có thứ tự và cây tìm kiếm nhị phân; cây nhị phân cân bằng; cây nhị phân
  tự điều chỉnh.

  \item[Chương 5. Nối và cắt cây]
  Bài toán nối và cắt cây; biểu diễn cây thành các tập đường; biểu diễn đường
  bằng cây nhị phân; ghi chú.

  \item[Chương 6. Cây khung nhỏ nhất]
  Phương pháp tham lam; ba thuật toán cổ điển; thuật toán vòng tròn; ghi chú.

  \item[Chương 7. Đường đi ngắn nhất]
  Cây đường đi ngắn nhất, gán nhãn và quét; thứ tự quét hiệu quả; mọi cặp
  đỉnh.

  \item[Chương 8. Luồng mạng]
  Luồng, lát cắt và đường tăng; tăng theo luồng chặn; tìm luồng chặn; luồng
  chi phí nhỏ nhất.

  \item[Chương 9. Ghép cặp]
  Ghép cặp hai phía và luồng mạng; đường luân phiên; blossom; thuật toán cho
  ghép cặp không hai phía.
\end{description}

\section{Lời nói đầu}

Trong mười lăm năm trước khi cuốn sách này được viết, lĩnh vực thuật toán tổ
hợp đã phát triển bùng nổ. Dù phần lớn công trình mới có tính lý thuyết,
nhiều thuật toán vừa được khám phá lại rất thực dụng. Các thuật toán ấy không
chỉ dựa trên kết quả mới trong tổ hợp và đặc biệt là lý thuyết đồ thị, mà còn
dựa trên sự phát triển của cấu trúc dữ liệu mới và kỹ thuật phân tích thuật
toán mới.

Mục tiêu của tác giả là làm rõ sự tương tác giữa các hướng này bằng cách giải
thích những thuật toán hiệu quả nhất đã biết cho một số bài toán tổ hợp chọn
lọc. Sách bao quát bốn bài toán kinh điển trong tối ưu mạng, kèm theo việc xây
dựng các cấu trúc dữ liệu cần dùng và phân tích thời gian chạy. Tác giả dự
định đưa phần này vào một bộ sách hai tập toàn diện hơn về cấu trúc dữ liệu và
thuật toán đồ thị.

Ba mục tiêu được nêu rõ là chiều sâu, độ chính xác và sự đơn giản. Tác giả cố
gắng trình bày các kỹ thuật tiên tiến nhất lúc bấy giờ theo cách dễ hiểu và
có thể dùng trong thực tế. Điều ông muốn truyền đạt là vẻ đẹp và chiều sâu của
thuật toán đồ thị, kiến thức về các thuật toán tốt nhất cho những bài toán
được xét, và cách hiện thực chúng.

Nội dung sách dựa trên các bài giảng tại hội nghị CBMS Regional Conference ở
Worcester Polytechnic Institute tháng 6 năm 1981. Sách cũng bao gồm công trình
chưa công bố khi đó của Tarjan cùng Daniel Sleator tại Bell Laboratories. Tác
giả cảm ơn ban tổ chức hội nghị, các học viên tham gia, National Science
Foundation, những người chuẩn bị bản thảo, Michael Garey vì các nhận xét sắc
sảo, và đặc biệt là Dan Sleator.

\section{Chương 1: Nền tảng}

Chương đầu đặt khung cho toàn bộ sách. Tarjan xét các thuật toán hiệu quả cho
bốn bài toán tối ưu mạng kinh điển. Những thuật toán này kết hợp hai mảng:
cấu trúc dữ liệu và phân tích thuật toán ở một bên, tối ưu mạng ở bên kia.
Tối ưu mạng lại liên quan tới vận trù học, khoa học máy tính và lý thuyết đồ
thị.

Một điểm nhấn quan trọng là: nhiều trình bày về thuật toán mạng bỏ qua chi
tiết cấu trúc dữ liệu, khiến người muốn cài đặt thuật toán phải tự giải một
bài lập trình không nhỏ. Tarjan chọn cách trình bày ở mức các thao tác đơn
giản trên những đối tượng toán học nguyên thủy như danh sách, cây và đồ thị,
thay vì viết chương trình cụ thể bằng FORTRAN.

Chương này cũng nhắc tới vai trò của độ phức tạp tính toán như một dạng dao
cạo Occam. Thuật toán hiệu quả nhất thường là thuật toán chỉ tính đúng lượng
thông tin cần thiết cho tình huống bài toán. Vì vậy, phát triển một thuật toán
đặc biệt hiệu quả thường đem lại hiểu biết sâu hơn về chính bài toán; kết quả
không chỉ nhanh mà còn đơn giản và đẹp.

\section{Chương 2: Các tập rời nhau}

Bài toán tập rời nhau yêu cầu duy trì một phân hoạch của tập phần tử dưới các
thao tác tìm đại diện và hợp hai tập. Cấu trúc cây nén, thường được biết đến
qua \emph{union-find}, là ví dụ kinh điển cho phân tích khấu hao: một thao tác
có thể tốn nhiều thời gian, nhưng trung bình trên cả dãy thao tác thì rất nhỏ.

Hai ý tưởng cơ bản là hợp theo hạng/kích thước và nén đường đi. Khi kết hợp,
chúng cho thời gian gần tuyến tính, với hệ số nghịch đảo Ackermann rất chậm
tăng. Đây là cấu trúc nền cho cây khung nhỏ nhất và nhiều thuật toán đồ thị
khác.

\section{Chương 3: Heap}

Heap là cấu trúc dữ liệu cho hàng đợi ưu tiên. Sách bàn về cây có thứ tự heap,
các biến thể đơn giản và heap trái. Trong tối ưu mạng, hàng đợi ưu tiên xuất
hiện tự nhiên trong Prim, Dijkstra và nhiều thuật toán gán nhãn.

Điểm quan trọng không chỉ là biết một heap chuẩn, mà là chọn biến thể phù hợp
với kiểu thao tác cần hỗ trợ: chèn, lấy phần tử nhỏ nhất, giảm khóa, và ghép
hai heap. Heap trái là một ví dụ cho việc thay đổi hình dạng cây để thao tác
ghép được nhanh.

\section{Chương 4: Cây tìm kiếm}

Chương này nhắc lại cây tìm kiếm nhị phân cho tập có thứ tự, rồi chuyển sang
cây cân bằng và cây tự điều chỉnh. Cây cân bằng giữ chiều cao \(\Oh(\log n)\)
để bảo đảm tìm kiếm, chèn và xóa hiệu quả. Cây tự điều chỉnh, tiêu biểu là
splay tree của Sleator và Tarjan, không lưu thông tin cân bằng tường minh mà
dựa vào các phép xoay sau mỗi truy cập.

Ý nghĩa trong sách là chuẩn bị cho những cấu trúc động phức tạp hơn, đặc biệt
là biểu diễn đường bằng cây nhị phân ở chương kế tiếp.

\section{Chương 5: Nối và cắt cây}

Bài toán \emph{link-cut tree} yêu cầu duy trì một rừng động dưới các thao tác
nối hai cây, cắt một cạnh, và truy vấn/cập nhật trên đường trong cây. Tarjan
trình bày ý tưởng biểu diễn mỗi cây bằng một tập các đường, rồi biểu diễn mỗi
đường bằng một cây nhị phân phụ.

Cấu trúc này là một trong các đóng góp nổi tiếng của Sleator--Tarjan. Nó cho
phép đưa nhiều thuật toán đồ thị động hoặc thuật toán tối ưu mạng từ mức ý
tưởng xuống mức cài đặt hiệu quả.

\section{Chương 6: Cây khung nhỏ nhất}

Chương này bắt đầu phần thuật toán mạng. Bài toán cây khung nhỏ nhất là môi
trường tự nhiên để thấy phương pháp tham lam hoạt động chính xác. Các thuật
toán cổ điển gồm Kruskal, Prim và Boruvka. Chúng khác nhau chủ yếu ở cách chọn
cạnh an toàn tiếp theo và cấu trúc dữ liệu hỗ trợ.

Union-find làm Kruskal hiệu quả; heap hỗ trợ Prim; còn các biến thể như thuật
toán vòng tròn khai thác cách tổ chức cạnh và thành phần liên thông để cải
thiện cận thời gian trong một số mô hình.

\section{Chương 7: Đường đi ngắn nhất}

Chương này trình bày cây đường đi ngắn nhất, thuật toán gán nhãn và thao tác
quét. Các thuật toán kinh điển như Dijkstra và Bellman--Ford có thể được nhìn
chung qua cách chọn thứ tự quét đỉnh/cạnh. Khi trọng số không âm, hàng đợi ưu
tiên cho phép chọn nhãn nhỏ nhất tiếp theo; khi có trọng số âm nhưng không có
chu trình âm, cần quét lặp để lan truyền cải thiện.

Phần mọi cặp đỉnh nối bài toán đơn nguồn với các phương pháp quy hoạch động
và nhân ma trận theo đại số min-plus.

\section{Chương 8: Luồng mạng}

Luồng mạng được trình bày qua định nghĩa luồng, lát cắt và đường tăng. Từ
định lý max-flow min-cut, các thuật toán có thể được xây bằng cách liên tục
tìm đường tăng trong mạng dư.

Điểm nâng cấp là dùng \emph{blocking flow}: thay vì tăng một đường mỗi lần,
ta xây đồ thị mức và tìm một luồng chặn để loại bỏ cả một lớp đường tăng
ngắn nhất. Sách cũng bàn cách tìm luồng chặn và chuyển sang luồng chi phí nhỏ
nhất, nơi ngoài giá trị luồng còn cần tối ưu tổng chi phí.

\section{Chương 9: Ghép cặp}

Chương cuối xét ghép cặp. Với đồ thị hai phía, ghép cặp cực đại có thể quy về
luồng mạng hoặc xử lý bằng đường tăng luân phiên. Với đồ thị không hai phía,
khó khăn chính là chu trình lẻ. Thuật toán blossom của Edmonds co các chu
trình lẻ này thành siêu đỉnh để tiếp tục tìm đường tăng.

Sách nhấn mạnh cùng một mẫu: hiểu cấu trúc tổ hợp của bài toán, rồi chọn cấu
trúc dữ liệu đủ mạnh để duy trì đúng thông tin cần thiết. Ghép cặp không hai
phía là ví dụ rõ ràng cho việc một trở ngại lý thuyết, blossom, phải được
biểu diễn cẩn thận trong thuật toán.

\section{Ghi chú sử dụng}

Cuốn sách này không phải sổ tay cài đặt từng dòng. Giá trị chính của nó là
cách nhìn thống nhất: các thuật toán mạng nhanh không tách rời cấu trúc dữ
liệu. Khi đọc trong bối cảnh thi thuật toán, nên đặc biệt chú ý các cặp sau:
\begin{itemize}
  \item Kruskal đi cùng union-find;
  \item Prim và Dijkstra đi cùng hàng đợi ưu tiên;
  \item cây động đi cùng link-cut tree;
  \item luồng chặn đi cùng đồ thị mức và cấu trúc quét cạnh;
  \item ghép cặp tổng quát đi cùng blossom và đường luân phiên.
\end{itemize}

\section*{Ghi chú về bản dịch}

PDF gốc có lớp văn bản tiếng Anh đọc được. Vì đây là một chuyên khảo dài, bản
LaTeX hiện tại chuyển ngữ phần trước sách và tạo hướng dẫn chương bằng tiếng
Việt, giữ các thuật ngữ thuật toán quan trọng để phục vụ tra cứu trong kho OI
Public Library.

\end{document}
